\question
Участники студенческого хакатона ИТМО отмечают каким языком
программирования они владеют. Пусть $U = \{1, \ldots, 8\}$~---
множество участников, $G \subseteq U$~--- знают Golang,
$C \subseteq U$~--- знают C++.

\[
    G = \{1, 2, 3, 4, 5, 6\}, \quad
    C = \{3, 4, 5, 6, 7, 8\}, \quad
    R = \{(1, 3), (3, 1), (4, 6), (6, 4), (5, 6), (6, 5), (5, 8), (8, 5)\}.
\]

\begin{enumerate}
    \renewcommand{\labelenumi}{\arabic{enumi}.}
    \item Найдите $M = G \cap C$~--- участников, знающих оба языка.

    \item Дано отношение $R$~--- пара участников может работать очно
    (симметричное). Рассматриваем только контакты между участниками,
    попавшими в множество $M$ (п.~1). Найдите пары, связанные очным
    контактом ровно через одного посредника. Обозначьте через $N$
    множество всех участников, которые вошли хотя бы в одну такую пару.

    \item Для найденного в пункте~2 множества $N$ вычислите
    $\overline{N}$ (дополнение относительно $U$) и $|P(N)|$.
\end{enumerate}

---------------

Автор -- Овчинников Михаил Денисович
